\section{Discussion}
\label{sec:discussion}

\paragraph{Convergence on inputs, not on mechanisms.} Seven forecasters with different tokenizers, depths and training
objectives single out the same series for a concept, share whole-layer geometry, and contain features with the same
broad effects: raising or lowering the forecast's level and damping or amplifying its seasonality. At the grain of
individual causal effects the convergence stops. Fine concepts are organized by model, the same capability can be one
head or a five-head circuit, and the same effect on the same series is detected only within one model family. One
reading consistent with all of this is that the models converge on \emph{what} in a series is predictive, which the
data dictate, while the \emph{implementation} of the response is left to each architecture. This is a weaker form of
convergence than the shared representations argued for in vision and language \citep{li2016convergent,huh2024platonic},
and our instruments locate where it stops: agreement on inputs is measurable and replicates, whereas agreement on effects
across architecture families is below what single-feature ablation resolves at these effect sizes.

\paragraph{What this means for interpreting TSFMs.} First, a single dictionary feature in a TSFM is only weakly special
against directions the model actually uses, so single-feature claims should be stated relative to an on-manifold null,
and the field's growing body of single-model SAE results
\citep{mishra2026dissecting,oublal2026timesae,divo2026zeitgeist} should be read with that in mind. Second, which
features matter must be measured by intervention: activation prominence, attention mass and correlation with a
generative factor each point at the wrong components. Third, cross-model comparison needs a measured time axis and a
pooling that keeps within-window order; otherwise similarity is understated by up to 0.43 CKA. Fourth, exploratory
pipelines with thousands of comparisons \citep{gelman2014garden} benefit from registration and a one-shot held-out test,
which here confirmed between half and two thirds of the registered claims (43 of 82 and 41 of 63) and sorted the
claim types cleanly into those that replicate and those that do not.

\section{Limitations}
\label{sec:limitations}

\paragraph{Models and alignment.} The study covers seven transformer models with contiguous tokens; models whose tokens
read disjoint lags, such as Lag-Llama, are out of scope. With seven models, differences between models are confounded
with architecture, size and training data, and Timer's features do not exceed chance, so its results carry no evidence
about causal concepts. The analysis is univariate (Chronos-2's cross-series attention is not measured; \fid{PM-13}), and
Chronos-T5's decoder and Sundial's flow head are not captured. The common time grid is no finer than the coarsest token,
and peak layer locations, unlike the size of similarity, depend on the window size (\fid{SH-23}). Dictionaries are
trained on window vectors, which is conservative for one-step models (\fid{MP-11}).

\paragraph{Ground truth and scope.} Ground truth is exact only for the synthetic tier and encodes the generators'
assumptions, and overlap with the models' pretraining data cannot be verified. Raw real benchmarks are a weaker basis for
causal questions, because weakly predictable real series receive flat forecasts in 67--100\% of cases (\fid{BM-06};
cf.\ \citealp{wan2026collapse}), leaving an ablation nothing to move; real data therefore enter as real-derived series
and as the external test set on which the transfer claims confirm. The confirmed transfers derive from two source
concepts on v2 and three source models on v3, so the shared-input result rests on few concepts. Concept results depend
on dictionary size, sparsity and seed; three replicate seeds measure this dependence but cannot remove it.

\section{Conclusion}
\label{sec:conclusion}

Across seven time-series foundation models and two independent corpora, with every causal instrument validated on a
planted known answer and every registered claim tested once on sealed data, architectures that look nothing alike agree
on what in a series matters---they share representational geometry, select the same inputs for a concept and contain
features that raise or lower the forecast's level and damp its seasonality---but implement the response in
model-specific, distributed ways. Several habits carried over from language-model interpretability, from averaging
tokens to trusting prominent features and attention maps, give confidently wrong answers on these models. The
benchmark, the analysis system and every result with its reproduction command are public, so each conclusion can be
re-tested on new models as they appear.

\paragraph{Availability.} Code (MIT license), configurations, the example run and the findings ledger are at
\url{\repourl}. Version 1.1.0 is archived at \href{https://doi.org/10.5281/zenodo.23267462}{doi:10.5281/zenodo.23267462},
and both packages install from PyPI (\texttt{pip install tsfm-lens tsfm-benchmark}).
